\chapter{Симулациски профили, mapper и безбедносен мост}
\label{chap:simulation}

\section{Интегрираниот ROS тек}

Симулацискиот пат е прикажан на \cref{fig:simulation-flow}. Mapper-от работи на
20 Hz, применува low-pass филтри за valence и arousal, quintic interpolation и
slew limits. Потоа safety bridge-от повторно ги ограничува командите,
арбитрира со manual input, проверува readiness/health и го објавува
единствениот \code{sensor\_msgs/Joy} transport на
\code{/emotion\_bot/joy\_out}. \code{/emotion\_bot/safe\_cmd} е inspectable
\code{Twist}, а не директен actuator transport.

\begin{figure}[htbp]
\centering
\resizebox{\textwidth}{!}{\begin{tikzpicture}[node distance=9mm and 10mm]
  \node[process,text width=25mm] (state) {state 1.1\\$v$, $a$, category};
  \node[process,right=of state,text width=29mm] (mapper) {mapper 20 Hz\\filter + keyframes};
  \node[safety,right=of mapper,text width=30mm] (safe) {safety bridge\\clamp + watchdog};
  \node[process,right=of safe,text width=25mm] (joy) {\code{/joy\_out}\\Joy};
  \node[process,right=of joy,text width=27mm] (ctrl) {Lite3 controller\\Gazebo};
  \draw[flow] (state) -- (mapper);
  \draw[flow] (mapper) -- node[above,font=\scriptsize]{Twist/action} (safe);
  \draw[flow] (safe) -- (joy);
  \draw[flow] (joy) -- (ctrl);
  \node[process,below=of safe,text width=29mm] (manual) {manual Joy input\\0.75 s priority};
  \draw[flow] (manual.north) -- node[right,font=\scriptsize]{arbitration} (safe.south);

  \node[danger,above=of safe,text width=32mm] (health) {readiness + health\\freshness + bounds};
  \draw[flow] (health.south) -- node[right,font=\scriptsize]{gates} (safe.north);

  \node[process,above=of joy,text width=25mm] (status) {status JSON\\read-only};
  \draw[dashedflow] (safe.north east) to[out=35,in=205]
    node[below,sloped,font=\scriptsize,fill=white,inner sep=1pt]{publish}
    (status.south west);
\end{tikzpicture}
}
\caption{Симулациски ROS тек од emotion-state до Lite3 controller.}
\label{fig:simulation-flow}
\end{figure}

Integrated launch-от стартува Gazebo, модел, 13 controllers, controller loader,
Lite3 runner, chat/emotion adapters, mapper и safety bridge. Controller
readiness се објавува на \code{/emotion\_bot/sim\_controller\_ready}. За
interactive handoff, status мора да покаже
\code{sim\_ready=true}, \code{health\_ok=true},
\code{stance\_prepared=true} и \code{motion\_enabled=true}.

\section{Fluid mapper}

Секој профил во \code{config/default.yaml} има entrance \code{segments} и
повторливи \code{idle\_segments}. Keyframe endpoints се поврзуваат со quintic
smootherstep

\begin{equation}
  s(u)=6u^5-15u^4+10u^3, \qquad u\in[0,1],
  \label{eq:smootherstep}
\end{equation}

за кој $s'(0)=s'(1)=s''(0)=s''(1)=0$. Така position, velocity и acceleration
немаат дисконтинуитет на seam. Mapper-от дополнително користи valence/arousal
time constants од 0.45/0.35 s, 1.25 entrance gain, 65\% idle amplitude и 1.25
idle duration factor.

Секоја non-neutral промена прво го откажува старото action generation, се
враќа на точна нула за 0.25 s и ја задржува canonical stance 0.35 s. Ова е
симулациски timing; физичкиот exact-neutral return е подолг, 1.5 s. Same-category
appraisal ја ажурира интензивноста без да го рестартира loop-от.

Normal profiles имаат

\begin{equation}
  v_x=v_y=\omega_z=0,
  \label{eq:centered}
\end{equation}

а ги користат body height, roll и pitch. Envelope-от е $\pm0.100$ m за height
и $\pm0.625$ rad за roll/pitch, со 0.18 m/s height slew и 1.00 rad/s mapper
attitude slew. Овие се симулациски граници и не се MotionSDK joint offsets.

\section{Девет профили}

Во \cref{tab:sim-emotions} се прикажани сите категории. Последните joint
excursions се од two-idle-cycle headless sweep на референтниот snapshot. Тие се
мерења на симулациски joints, а не препорачани hardware амплитуди.

\begin{table}[htbp]
\centering
\caption{Девет симулациски емоционални профили}
\label{tab:sim-emotions}
\small
\begin{tabularx}{\textwidth}{L{0.15\textwidth}YL{0.24\textwidth}}
\toprule
\textbf{Категорија} & \textbf{Визуелен израз во Gazebo} & \textbf{Последен измерен max idle joint excursion} \\
\midrule
Neutral & меко дишење и наизменичен tilt & 0.1081 rad \\
Joy & еден центриран hop, високи/ниски rocks и 1.40 s bounce loop & 0.2378 rad \\
Sadness & длабок bow и бавно ниско нишање & 0.1644 rad \\
Anger & симетричен multi-impact stomp и напнат низок loop & 0.4017 rad \\
Fear & брз recoil, crouch и страничен tremble & 0.1884 rad \\
Surprise & центриран hop, висок alert став и краток loop & 0.2932 rad \\
Disgust & асиметричен down-and-away recoil & 0.2525 rad \\
Curiosity & широки наизменични body tilts & 0.1026 rad \\
Affection & бавно, топло sway движење & 0.1104 rad \\
\bottomrule
\end{tabularx}
\end{table}


Joy и Surprise користат controller-side hop од 0.55 s со вертикални offsets
\begin{equation*}
(-0.069,\ +0.100,\ -0.050,\ 0.0)\ \mathrm{m}.
\end{equation*}
Anger користи 0.75 s stomp со седум offsets и
smooth recovery. Двете actions се bounded, cancellation-aware и generation
ordered. Максималната вертикална брзина е 1.125 m/s. Не се применува Gazebo
external wrench. Поради тоа „hop“ и „stomp“ во ова поглавје значат исклучиво
симулациски controller actions; тие не се физички airborne jump или torque
strike.

\section{Safety bridge, arbitration и watchdog}

Состојбите од \cref{fig:simulation-safety} ја илустрираат operational logic.
Системот започнува со нулта команда. По controller readiness се подготвува
\code{JOY\_STAND} и се чека четири секунди simulated time за settling. Manual
activity има приоритет 0.75 s, но ги дели истите clamps. Emotion и manual
командите истекуваат по 0.5 s; emotion state по 1.0 s.

\begin{figure}[htbp]
\centering
\begin{tikzpicture}[node distance=17mm and 22mm]
  \node[box,text width=27mm] (zero) {\textbf{ZERO}\\exact zero};
  \node[process,right=of zero,text width=27mm] (ready) {\textbf{READY}\\stance prepared};
  \node[safety,right=of ready,text width=27mm] (enabled) {\textbf{ENABLED}\\bounded output};
  \node[danger,below=of ready,text width=27mm] (fault) {\textbf{FAULT}\\permission revoked};
  \node[physical,below=of enabled,text width=27mm] (recover) {\textbf{RECOVERY}\\recentering};

  \draw[flow] (zero) -- node[above,font=\scriptsize]{controllers + stance} (ready);
  \draw[flow] (ready) -- node[above,font=\scriptsize]{enable} (enabled);
  \draw[flow] (enabled) -- node[right,font=\scriptsize]{soft drift} (recover);
  \draw[flow] (recover.west) to[bend left=22]
    node[below,font=\scriptsize,fill=white,inner sep=1pt]{inside radius} (enabled.south west);
  \coordinate (stalelevel) at ($(ready.north)+(0,8mm)$);
  \draw[dashedflow] (enabled.north) |- (stalelevel) -| (zero.north);
  \node[above=0.5mm,font=\scriptsize,fill=white,inner sep=1pt]
    at (stalelevel) {stale / disable};
  \draw[flow] (enabled.south west) --
    node[above,sloped,font=\scriptsize,fill=white,inner sep=1pt]{health / bound fault}
    (fault.north east);
  \draw[flow] (fault.west) -| node[pos=0.25,below,font=\scriptsize]{explicit reset} (zero.south);
\end{tikzpicture}

\caption{Поедноставена состојбена машина на симулацискиот safety bridge.}
\label{fig:simulation-safety}
\end{figure}

Заштитата следи finite model pose, torso height, tilt и 12 joints. Invalid data,
health gap, disable, stale source или node loss водат кон exact zero и revoked
permission. Симулаторот останува dynamic, но има simulation-only recenter:
на 0.09 m од enable-time центарот започнува recovery со најмногу 0.012 m
increment, сè додека роботот не влезе под 0.035 m; 0.20 m е hard stop. HipX
anti-splay term од 40 Nm/rad ги центрира четирите abduction joints. Ниту еден
од овие Gazebo механизми не е физички controller claim.

\section{Симулациски verification резултати}

На референтниот snapshot, \code{make -C lite3-noetic verify-emotion} е
repeatable offline authority. Датираниот record наведува:

\begin{itemize}
  \item clean Release \code{catkin\_make} со само постојни upstream warnings;
  \item 58 Python tests и 5 controller-side C++ trajectory tests;
  \item еден offline OpenAI process-boundary streaming test;
  \item еден ROS integration test со 0 failures/errors;
  \item два isolated headless Gazebo sweeps со 13 controllers и 12 joints;
  \item сите девет entrance/idle profiles, rapid Joy$\rightarrow$Anger$\rightarrow$Fear,
  manual priority, watchdog zeroing и no-hardware boundary; и
  \item per-profile contact-model drift под 0.025 m во најновиот two-loop sweep,
  при guard од 0.10 m.
\end{itemize}

Anger произвел три целосни stomp occurrences и recoveries; Joy и Surprise ги
извршиле hop actions. Секоја директна non-neutral промена имала најмалку пет
последователни exact-zero mapper samples пред нов entrance. WSLg GUI бил
прегледан со блиска камера, при што повеќе chat-driven transitions останале
исправени и со разделени стапала. Овие резултати се силен симулациски доказ,
но не докажуваат физичко однесување.
